% Pentest_report.tex
\documentclass[a4paper]{article}

\usepackage[utf8]{inputenc}
\usepackage[spanish]{babel}
\usepackage{graphicx} % Para la inserción de imagenes
\usepackage{lmodern}
\usepackage[dvipsnames,svgnames]{xcolor}
\usepackage[margin=2cm,top=2cm, includefoot]{geometry}
\usepackage[most]{tcolorbox}
\usepackage{fancyhdr}
\usepackage[hidelinks]{hyperref}
\usepackage{parskip}
\usepackage{listings}
\usepackage{smartdiagram}

% Declaración de colores
\definecolor{greenPortada}{RGB}{105,168,79}

% Declaración de variables
\newcommand{\logoPortada}{images/logo.png} % Update this path to match your logo location
\newcommand{\logoMaquina}{images/samba.png} % Update this path to match your machine logo location
\newcommand{\machineName}{Lame} % Nombre de la máquina

% Adicionales
% Ajuste de headheight recomendado por fancyhdr (evita el error "headheight is too small")
\setlength{\headheight}{60pt} % >= 54.67494pt sugerido por fancyhdr
\addtolength{\topmargin}{-14.47495pt} % opcional, sugerido por fancyhdr para mantener el diseño

\pagestyle{fancy}
\fancyhf{}
\lhead{\includegraphics[width=2cm]{\logoPortada}}\rhead{\includegraphics[width=1.2cm]{\logoMaquina}}
\renewcommand{\headrulewidth}{3pt}
\renewcommand{\headrule}{\hbox to\headwidth{%
    \color{greenPortada}\leaders\hrule height \headrulewidth\hfill}}

% Comienzo del documento
\begin{document}
        \cfoot{\thepage}
% Creacion de la portada
\begin{titlepage}
        \centering
        \includegraphics[width=0.6\textwidth]{\logoPortada}\par\vspace{1cm}
        {\LARGE \textbf{WRITE-UP}} \par\vspace{0.25cm}
        {\Huge\bfseries\textcolor{greenPortada}{Maquina \machineName}\par\vspace{1cm}}
        \vfill
        \includegraphics[width=\textwidth,height=6cm, keepaspectratio]{\logoMaquina}\par\vspace{0.5cm}
        \vfill
        \begin{tcolorbox}[colback=red!5!white,colframe=red!75!black]
        \centering
        {\Large Autor: Aarón Ibarra}\par\vspace{0.4cm}
        {\large Fecha: 10 de Noviembre de 2025}\par
        \end{tcolorbox}\par\vspace{0.1cm}
        \vfill
        {\large Todos los derechos reservados.}
        \vfill
        \end{titlepage}

        % Índice -------------------------------------------------------------------------------------
                \clearpage
        \tableofcontents
        \clearpage

        % Contenido del informe ----------------------------------------------------------------------
        % Introducción
        \section{Introducción}
        \vspace{0.2cm}

        \subsection{Antecedentes}
        La máquina {\textbf{\machineName}} de la plataforma \href{https://hackthebox.com}{\textbf{\textcolor{blue}{HackTheBox}}} pensada para principiantes: fue una de las primeras máquinas públicas y es ideal para aprender enumeración y explotación de servicios antiguos. Es una VM Linux con servicios como Samba (versión vulnerable y  exploit asociado a CVE-2007-2447). \vspace{0.2cm}
        \begin{figure}[h]
                \centering
                
\includegraphics[width=0.55\textwidth]{images/Lame details.png}
                \caption{Detalles de la máquina.}
        \end{figure}
        \vspace{0.2cm}
        \begin{tcolorbox}[enhanced,attach boxed title to top center={yshift=-3mm,yshifttext=-1mm},
                colback=blue!5!white,colframe=blue!75!black,colbacktitle=greenPortada!80!black,
                title=Dirección URL,fonttitle=\bfseries,
                boxed title style={size=small,colframe=blue!50!black} ]
                \centering
                \href{https://app.hackthebox.com/machines/lame}{\color{blue}{https://app.hackthebox.com/machines/lame}}
        \end{tcolorbox}
        \vspace{0.2cm}
        
        \subsection{Objetivos}
        Durante este ejercicio, se busca identificar y enumerar servicios vulnerables, descubrir y explotar vulnerabilidades que permitan obtener acceso remoto (foothold), y documentar hallazgos con pasos reproducibles.            
        \vspace{0.5cm}
        \begin{figure}[h]
                \centering
                \smartdiagram[priority descriptive diagram]
                {Reconocimiento para delimitar alcance y contexto operativo,
                Escaneo y enumeración de servicios y dependencias,
                Explotación controlada para verificar vulnerabilidades,
                Postexplotación segura y documentación de hallazgos}
                \caption{Fases del ataque.}
        \end{figure}
        \clearpage

        % Enumeración
        \section{Enumeración}
        \subsection{Reconocimiento inicial}
        Se inició un análisis preliminar para comprobar la accesibilidad del sistema objetivo desde la máquina atacante, empleando la herramienta \texttt{ping} para enviar paquetes ICMP y confirmar la conexión del equipo en la red.\vspace{0.2cm}
        \begin{figure}[h]
                \centering
                \includegraphics[width=0.6\textwidth]{images/ping.png}
                \caption{Reconocimiento inicial sobre el sistema objetivo.} 
        \end{figure}
        \vspace{0.2cm}

        Tras localizar el sistema se realizó un escaneo de puertos con la herramienta \texttt{nmap} con el fin de identificar los servicios en ejecución y las versiones asociadas, registrándose los siguientes hallazgos.\vspace{0.2cm}
        \begin{figure}[h]  
                \centering
                \includegraphics[width=.75\textwidth]{images/nmap-scan.png}
                \caption{Escaneo inicial de puertos con nmap.}  
        \end{figure}
        \vspace{0.2cm}
        \clearpage
        \subsection{Enumeración de servicios}
        A continuación, se procedió a una enumeración más detallada de los servicios identificados en el escaneo inicial, con el fin de obtener información adicional sobre sus configuraciones y posibles vulnerabilidades asociadas.\
        \vspace{0.2cm}
        \begin{figure}[h]
                \centering
                \includegraphics[width=.6\textwidth]{images/nmap-services.png}
                \caption{Enumeración detallada de los servicios con nmap.}
        \end{figure}
        \vspace{0.2cm}
        
        % Explotación
        \section{Explotación}
        Durante la etapa de reconocimiento se determinó que el host objetivo ejecutaba \textbf{Samba 3.0.20}. Con esta identificación se realizó una búsqueda dirigida de vulnerabilidades y exploits relevantes mediante \textbf{searchsploit} (Metasploit), con el fin de seleccionar vectores de explotación aplicables a la versión encontrada.\vspace{0.2cm}        \vspace{0.2cm}
        \begin{figure}[h]
                \centering
                \includegraphics[width=.9\textwidth]{images/search-samba.png}
                \caption{Búsqueda de exploits para Samba 3.0.20 con searchsploit.}
        \end{figure}
        \vspace{0.2cm}
        \clearpage

        Se seleccionó el exploit asociado a la vulnerabilidad \textbf{CVE-2007-2447} (Samba 'username map script' Remote Code Execution), que permite la ejecución remota de comandos mediante la manipulación del parámetro \texttt{username map script} en la configuración de Samba. A continuación, se procedió a configurar el exploit con los siguiente parametros.\vspace{0.2cm}
        \begin{figure}[h]
                \centering
                \includegraphics[width=.9\textwidth]{images/show-options.png}
                \caption{Configuración y ejecución del exploit en Metasploit.}
        \end{figure}
        \vspace{0.2cm}

        Finalmente, se logró obtener una shell remota en el sistema objetivo, confirmando el éxito de la explotación y permitiendo el acceso al entorno del host comprometido.
        \begin{figure}[h]
                \centering
                \includegraphics[width=.9\textwidth]{images/exploitation.png}
                \caption{Obtención de shell remota en el sistema objetivo.}
        \end{figure}
        \vspace{0.2cm}  
        \end{document}

